\documentclass{brochure}

\begin{document}

\chapter{标准所涉概念体系}

国家标准GB 18030—2022《信息技术　中文编码字符集》（后文简作“本标准”）是一份为中文建立编码字符集的强制性国家标准。所以要理解这份标准的内容、功能和作用，既要懂中文的概念体系，又要懂编码字符集的概念体系。

但是首先，中文（Chinese scripts）在中文信息处理领域里的定义就和平常的用法很不一样，是“中国各族人民使用或曾经使用过的文字的统称”，这是一个囊括了中国各少数民族现行文字和中国民族古文字的概念，而不是特指中国的国家通用语言文字的文字部分的概念。\sidenote{旁注：参考代红《GB 18030中文编码字符集标准实施要求探讨》，载《信息技术与标准化》，2020年第6期。}所以从标准的几个版本里，我们既可以看到规范汉字，又可以看到不断增添进来的少数民族现行文字和中国民族古文字。

所以这里说的要懂“中文”的概念体系，就不仅仅是懂规范汉字，或者懂汉字文字学的知识，还要能在一般的角度把握文字，而这主要是比较文字学和普通文字学领域的内容。比较文字学和普通文字学的建立比中文进入计算机要早很多，所以从一般的角度介绍文字也可以说是在介绍字符集出现以前的时代里的概念体系，也就是“前字符集时代”的概念体系。先有了对文字的一般把握，在理解编码字符集时代和字符集时代的概念体系的时候就会方便轻松一些。

\section{前字符集时代概念体系}

\subsection{作为书写系统的文字}

在面对各种各样的文字时，首先遇到的问题是称呼问题。\sidenote{“文字”既可以指一个个的文本元素，也可以用来指一整个体系。本文用“文字”时指的都是文字体系。}最常见的说法是在民族的名称后面加上“文”字，比如汉族人用汉文，蒙古族人用蒙古文，藏族人用藏文，傣族人用傣文，等等。但在考察景颇族的语言文字时，发现景颇族的景颇支系和载瓦支系分别使用景颇文和载瓦文，这同景颇语和载瓦语作为同一民族不同支系所操的不同语言也有关系。\sidenote{参考《论景颇语和载瓦语的关系》，载《戴庆厦文集·第一卷》，中央民族大学出版社2012年，原载《思想战线》1981年第4期。}所以在民族的名称后面加上“文”字，不如说是在语言的名称后面加上“文”字，比如记录汉语的是汉文，记录蒙古语的是蒙古文，记录藏语的是藏文，记录傣语的是傣文，等等.

进一步考察发现，即使是记录同一语言的文字，在不同历史时期也会使用不同样式的字形系统。比如蒙古族最初采用回鹘字母记录蒙古语，称回鹘蒙古文，随后发展为胡都木蒙古文和托忒蒙古文，20世纪50年代初还曾试用过西里尔蒙古文。又如维吾尔族使用阿拉伯字母记录维吾尔语，称阿拉伯维吾尔文，1964年推行使用拉丁字母记录维吾尔语，称拉丁维吾尔文。这些不同的字母体系，如回鹘字母、阿拉伯字母、西里尔字母、拉丁字母，可以称为不同的图符集。

进一步考察发现，即使是记录同一语言的使用同一套图符集的文字，在不同历史时期还会制订不同的规范体系。比如维吾尔族在使用阿拉伯字母记录维吾尔语后，不断对这套文字进行改进和补充，原有阿拉伯维吾尔文包含32个字母，到1937年改进的阿拉伯维吾尔文仍有32个字母，1954年制订的正字法规则（称1954年版阿拉伯维吾尔文）中包含30个字母，到1984年起推行的阿拉伯维吾尔文字母表（称1984年版阿拉伯维吾尔文）中则包含32个字母，字母顺序也有所调整。\sidenote{因为正字法的编制年份、正字法的审定年份和文字的推行年份可能不同，按不同的标准会形成不同的称呼。}又如壮族曾使用方块壮字记录壮语，后来使用拉丁字母记录壮语，其中1957年批准的拉丁壮文（称1957年版拉丁壮文）包含32个字母，1982年推行的拉丁壮文（称1982年版拉丁壮文）包含26个字母。

所以，在中文的范围内，要用三个部分才可以准确地称呼一种文字：首先是这种文字记录的语言（汉语、蒙古语、藏语、维吾尔语，等等）；然后是这种文字使用的图符集（拉丁字母、阿拉伯字母、西里尔字母，等等）；最后是这种文字使用的规范体系（拼音方案、构词规范、笔顺规范，也就是正字法、正词法、正形法，等等）。\sidenote{如果语言、图符集和规范体系因国家或行政区划而异，则还要说明是哪个国家或地区使用的才行。}用这三个部分来称呼的文字，也就是用于记录某种语言的、采用某一套或某一些图符集的、带有某一套规范的（有时是某一历史时期的某个用户社群使用的）文字，称为一套书写系统（writing system）。

\subsection{形素与形位}

书写系统既然是语言、图符集和规范体系关联在一起的产物，它的主要组成部分自然是语言和图符集的对应关系，比如字音（与字义）和字形的对应关系、音系和字母表的对应关系、音节和字表的对应关系，等等。语言学、音韵学以及语言信息处理领域倾向于研究对应关系的左侧（语言侧），而本文涉及的工作，包括比较文字学、普通文字学倾向于研究对应关系的右侧（文字侧），也就是偏重于研究和处理书写系统的图符集。所以首先要解决图符集从哪儿来的问题，《汉字构形学导论》给出了两种方案。\sidenote{参考王宁《汉字构形学导论》，商务印书馆2019年，162页。}

一种是从社会文本中整理出图符集。社会文本包括社会通行文本（官方正式文告、名人书写诗篇、有广泛影响力的典籍钞本，等等）与民间书写文本（个人书信、账目、便笺、日记、契约、影响力有限的典籍钞本，等等），是图符集在使用过程中的实现。通过分析社会文本和语言的对应关系，可以从社会文本中开展多层次的归纳。其中最基本的层次是从形素（graph）中归纳出形位（grapheme）。比如对于使用汉字的文本，我们可以将许多手写的“汉”字归纳成一个“汉”字，每个手写的“汉”字都是一个形素，归纳成的一个“汉”字是一个形位。\sidenote{王宁《汉字构形学导论》（123页）使用“从许多个字中归纳出字样”“对于不同书法家所写的楷书汉字，那些写法完全一样的，都属于同一个字样”的说法。但要注意，这里的“字样”比“形位”要严格，还要求记词功能相同、构形构意相同，因此只能说形位归纳和字样归纳都有在同一书写系统内沟通写法的作用。}

另一种是从规范文本中提取出图符集。规范文本包括正字法、正词法、正形法等，是直接规定书写系统的图符集的材料。规范文本的内容本身就有框定字样、廓清边界、拟定主形、限定文本组合排列结构、规定变形规则等的作用，所以解读规范文件就能以主形划定形位，并获知形位之间的关系。

\subsection{形位系统}

图符集就是形位的集合，也就是形位集（repertoire）。在文字规范的要求下，形位集中的形位之间可能有着特定的堆叠关系、排列关系、组合关系、关联关系，这使得形位集构成一个系统，称为形位系统（grapheme system）。

如果一个形位叫一个“字”，对应的形位集就叫“字汇”，按照书写系统给字分配的功能可以大致分出音节字汇（syllabic repertoire）、表意字汇（ideographic repertoire），等等。如果一个形位叫一个“字母”，对应的形位集就叫“字母表”。按照书写系统给字母分配的功能和字母之间的分工可以大致划分出希腊式字母表（alphabet）、印度式字母表（abugida）、阿拉伯式字母表（abjad），等等。\sidenote{这种划分是纯粹针对字母的功能来说的，和用不用阿拉伯字母、用不用希腊字母没有关系。比如1984年版阿拉伯维吾尔文用的是阿拉伯字母，但按照字母的功能，它应该分析成一套希腊式字母表。}

\section{编码字符集时代概念体系}

\subsection{字图}

“文字进入计算机”，说的是要让计算机能处理文字信息。但一般的计算机只能处理二进制位（比特，bit）组成的数据，所以问题转换成如何为文字信息来设计比特组成的数据，以及反过来的面对比特序列应该如何解析成文字信息。

对于字或者字母来说，它的图像是最方便用二进制位表示的。二进制位代表着两种状态，如果屏幕上的一个小格子亮的状态代表“1”，不亮的状态代表“0”，那么一个字或者字母对应的图像就是这个范围里屏幕亮和不亮的排成矩阵的二进制位信息，也叫作这个字或者字母的点阵字图（bitmap glyph）。\sidenote{后文都用“0”和“1”表示二进制位的两个状态。}“点阵”是点的阵列，表示这个图像由图像点的阵列构成，是个点阵图；“字图”是字或者字母的图像，表示这个图是对应这个字或者字母的。\sidenote{glyph是这个术语的英文简称，英文全称是glyph image，中文翻译成字图。一些标准用glyph来指代字形，本文改用form来指代字形。}

用点阵图来存储汉字的字形是非常方便的，因为汉字是方块字，每个字在特定的字号下面用到的二进制位是一样多的，于是所有的点阵图存储在一个文件里的时候，就能很快计算出第几个字是从这个文件的第几个二进制位开始的。这种把字图存储在一起的文件称为字图集文件（glyph file），以前也叫字模集文件，后来叫作字库文件（font file）。所以集中存储点阵图的文件叫点阵字库文件（bitmap font file）。

但是因为点阵图是平面图，当字号扩大，网格和点阵扩大，文件的大小是平方增长的。这导致大字号字库文件的大小容易超出内存的限制。为了压缩点阵字库文件的大小，王选提出了存储折线的轮廓信息和汉字笔画信息的数据结构。在字库文件中只用存储折线首尾的点的横纵坐标，而在渲染文字的时候再把折线转换成点阵。这样就不是在存储点阵图，而是折线轮廓图，对应的字库可以叫折线字库。后来又有使用光滑的曲线来存储字图，也就是存储曲线轮廓图，对应的字库可以叫曲线字库。

\end{document}